%-------------------------------------------------------------------------------
% Experience — DevRel variant (leadership-forward)
% Hello Soda, Consultant and USC entries are shared: edit them in
% cv-sections/shared/.
%-------------------------------------------------------------------------------
\cvsection{Experience}

\begin{cventries}

\cventry
{CEO of Impact}
{Quansight PBC}
{Manchester, UK}
{May 2025 -- Present}
{
\begin{cvitems}
  \item {Sustain and grow long-running partnerships that fund maintainer time on community-wide work, such as free-threaded Python with Meta; documentation, tooling and demos across the JAX and OpenXLA ecosystems with Google, including support for the Google Developer Experts community; and a multi-stakeholder effort with NVIDIA, Red Hat and Meta to improve packaging across the Python ecosystem.}
  \item {Brought in new partners including Red Hat, DeepMind and the Gates Foundation.}
  \item {Decide where Quansight invests across its open source portfolio (conda, conda-forge, Jupyter, Nebari and related projects), weighing what clients, maintainers, users and our own staff need most.}
  \item {Maintain, design, write and review tooling used by millions of developers and researchers every day; currently improving conda's environment resolution and package verification.}
  \item {Answer to the board for both sides of Quansight's mission as a Public Benefit Corporation: a sustainable business and a healthier open source ecosystem.}
\end{cvitems}
}

\cventry
{Director}
{Quansight Labs}
{Manchester, UK}
{January 2021 -- May 2025}
{
\begin{cvitems}
  \item {Grew Quansight Labs' open source work 200\% over the role, measured across staff hours, pull requests, time spent on community work such as calls and mentoring, and the number of projects we supported; aligned it with our consultancy work, growing revenue 25\%.}
  \item {Took over Nebari, an open source platform for deploying AI and ML infrastructure, and turned it from an internal tool into a community project with its own governance. As product owner for around two years, set the roadmap, fed community and client feedback into it, and wrote and reviewed code; users include NOAA and NASA.}
  \item {Created and led an open source mentorship programme, now in its sixth year, with 50 mentees across 20+ open source projects.}
  \item {Led a 50+ person team of software engineers, data scientists, ML engineers, UX designers and developer-experience engineers, setting technical direction across concurrent client and open source engagements.}
  \item {Designed and tested ways to keep open source sustainable (consortia, open core, collaborative funding, community-driven models) and advised other organisations on open source strategy and governance.}
  \item {Architected and led a computer-vision ML platform, built from scratch, that let a social marketing client detect brands across sports, esports and live-event video. It scaled to 10,000+ brands and 8+ hours of video a day, at 40\% lower cost than the product it replaced.}
\end{cvitems}
}

\cventry
{Sr. Developer Advocate \& Tech Lead}
{Microsoft}
{Manchester, UK}
{March 2019 -- December 2020}
{
\begin{cvitems}
  \item {Co-developed Microsoft's Python advocacy strategy for the education, scientific computing and data science communities: where to show up, what to build, and which partnerships to pursue.}
  \item {Brought feedback from the scientific and numerical Python communities into VS Code, working closely with the Jupyter notebook and Pylance teams on better plotting in notebooks (Matplotlib, Bokeh and other interactive plots), Python environment discovery, the Gather extension, and many smaller usability fixes.}
  \item {Worked with Azure leadership on competitive analysis of developer tools, shaping improvements to the Azure Functions CLI and SDK experience and to its offerings for different tiers of developers.}
  \item {Took research software needs to the Azure teams, and contributed features and performance improvements to VS Code, Azure ML, Azure Functions and the Azure SDK.}
  \item {Created tutorials, blog posts, videos and live streams, and delivered 120+ talks, workshops and keynotes, training 2,000+ developers, scientists and students.}
  \item {Ran meetups, sprints and hackathons for the Python and data science communities.}
  \item {Led cross-functional teams of developers, researchers and project managers on research partnerships, including a secure data platform for hundreds of data scientists at the Alan Turing Institute and a cloud platform for the UK's national weather service processing millions of observations a day.}
\end{cvitems}
}

%-------------------------------------------------------------------------------
% Hello Soda — shared experience entry, \input inside each variant's cventries
%-------------------------------------------------------------------------------
\cventry
{Sr. Machine Learning Engineer}
{Hello Soda}
{Manchester, UK}
{July 2018 -- February 2019}
{
\begin{cvitems}
  \item {Architected a platform running 60+ ML and NLP models for customers worldwide, improving model uptime and performance by 15\%.}
  \item {Introduced MLOps practices and automation that cut model deployment from 14 days to 2 and improved system reliability by 40\%.}
  \item {Put data governance and deployment standards in place that cut operational costs 25\%, improved model accuracy 20\% and reduced production incidents.}
\end{cvitems}
}

%-------------------------------------------------------------------------------
% Independent consultant — shared experience entry
%-------------------------------------------------------------------------------
\cventry
{Independent Consultant}
{Freelance}
{Manchester, UK}
{November 2016 -- Present}
{
\begin{cvitems}
  \item {Advise start-ups, researchers and organisations on reproducible data workflows, data pipelines from ETL to deployment, and data and open source governance, reducing computation time on data-heavy workflows in most engagements.}
\end{cvitems}
}


\cventry
{Research Software Engineer, ML Team}
{The University of Sheffield}
{Sheffield, UK}
{November 2016 -- June 2018}
{
\begin{cvitems}
  \item {Core contributor to OpenDreamKit (EU-funded) and Project Jupyter; built analysis pipelines in Python and R for NLP, bioinformatics, materials science and geospatial research.}
  \item {Mentored colleagues and students in reproducible workflows, software engineering and HPC.}
  \item {Served on the UK RSE and NumFOCUS diversity and inclusion committees, contributing to national and international policy.}
  \item {Built dashboards and visualisations with d3.js, Shiny, Bokeh and Plotly so researchers could share their findings.}
\end{cvitems}
}

%-------------------------------------------------------------------------------
% University of Southern California — shared experience entry
%-------------------------------------------------------------------------------
\cventry
{Research Engineer, Medical Devices}
{University of Southern California}
{California, USA}
{2011 -- 2012}
{}


\end{cventries}
